\frametitle{Starting Point, and the Question}
\footnotesize

\begin{itemize}
\item The starting point is a \textbf{classical, working solver}: the Laplace--Beltrami
problem with a variable coefficient, discretised on the surface of a torus with high-order
surface finite elements, solved by CG preconditioned by a geometric multigrid V-cycle.
\item Both parts are implemented and verified. The discretisation converges at the optimal
rates for the degree used, and the preconditioned iteration is \alert{mesh-independent} ---
it reduces the residual by $10^{-10}$ in essentially the same number of steps whether the
system has one hundred or ten thousand unknowns.
\end{itemize}

\vspace{1mm}
A learned smoother is therefore \emph{not} asked to repair a solver that fails. It cannot
distinguish itself by restoring scalability, because scalability is already there. What it
can be asked to do is narrower: reduce the \emph{constant}.

\begin{block}{The question this work answers}
\itshape
Not whether a network can solve the surface problem, and not whether it can replace the
multigrid solver. The question is whether a network can supply the correction that
\emph{one node} of the existing cycle computes.
\end{block}

\begin{itemize}
\item The \textbf{smoothing node} is the natural place to ask this, for a structural
reason rather than an opportunistic one: of smoothing, residual evaluation, restriction,
the coarse solve and prolongation, the relaxation sweep is the \emph{only} operation whose
replacement leaves the multilevel structure intact.
\item Everything else keeps its classical form, so a change in the cycle's behaviour is
\alert{attributable to the substituted node} rather than to a re-engineering of the method
--- and that literature has been shown to suffer from weak baselines and reporting bias
\cite{mcgreivy2024weak}.
\end{itemize}

\vspace{1mm}
Three stages, named from the beginning so that a Stage~I result is not read as more than it
is: \textbf{I} --- are the components the coarse grid cannot reach reduced? \emph{(measured;
this talk.)} \textbf{II} --- what does the step do to the ones it \emph{does} reach?
\emph{(designed.)} \textbf{III} --- does it survive the V-cycle? \emph{(protocol fixed in
advance.)}
